\section{Introduction}

Two independent runs of the same AI puzzle hunt pipeline, fed identical inputs --- the same Age of Empires~II world data, the same scope document, the same toolkit specification --- produced hunts that share zero answer words (Jaccard similarity: 0.00), use categorically different meta mechanisms (a crossword grid trace versus a vowel-count index followed by anagram resolution), adopt different round architectures (a single flat round versus a three-round gated structure), and arrive at different meta answers (\textsc{wololo} versus \textsc{grail}). At the same time, both hunts maintain identical narrator voice, identical audience specification, identical format, identical total puzzle count, and similar solve-time estimates. Both pass the same quality threshold under AI expert panel review with normalized scores within four percentage points of each other. The divergence is real, total, and systematic at some levels of the pipeline; the convergence is equally real, complete, and systematic at others.

This pattern poses a practical problem for practitioners deploying AI creative generators. Current pipelines for AI-assisted creative production --- whether for puzzle hunts, tabletop campaign design, interactive fiction, or procedural game content --- yield outputs that vary stochastically across runs. Practitioners face a binary choice in response: review every output decision manually (expensive, and defeats the purpose of automation) or trust the pipeline's outputs as-is (cheap, but potentially arbitrary where it should not be). Neither strategy is principled, because neither is informed by a characterization of which decisions are actually stable across runs and which are genuinely variable. Without this characterization, practitioners cannot know which outputs to run once and trust versus which to run multiple times and select from. They either over-invest human oversight where the AI is reliable, or under-invest where reliability is absent. A principled account of divergence structure is the necessary foundation for rational deployment of AI creative pipelines.

We address this problem through a \emph{same-input divergence} experiment: two fully independent runs of a complete puzzle hunt generation pipeline from identical input documents, compared at five levels of abstraction --- structural decisions, puzzle pool selection, answer word assignment, extraction mechanism design, and micro-level clue formulation. The experiment treats stochasticity not as a confound to be eliminated (as in reproducibility research) but as an instrument for revealing the structure of the pipeline's design space. Decisions where the two runs converge, despite stochastic generation, are decisions where the design space is constrained enough that independent generative agents reach similar conclusions. Decisions where the two runs diverge are decisions where the design space offers many equally valid options and nothing in the specification privileges one over another. The pattern of convergence and divergence, read across five levels, characterizes which parts of the pipeline are reliable and which are not.

Our central empirical finding is that divergence in this pipeline is structured, not random, and that the structure is predicted by a single principled distinction: \emph{specification drives convergence}. Every decision that was explicitly specified in the creative brief --- scale (five feeders plus one meta), narrator voice (the Monk persona), output format (printable Markdown), target audience (solo AoE2 player, casual), and solve-time target --- is stable across runs to the point of identity. Every decision left to the pipeline --- round count, round themes, meta mechanism family, answer word selection, extraction logic within puzzle types, clue formulations --- diverges substantially or totally. Notably, this is not a domain-knowledge effect: the Age of Empires~II domain is rich and well-defined, but domain richness neither constrains the round structure nor determines the meta answer. Domain content acts instead as what we call \emph{semantic gravity} at the word-selection level: even when answer word sets are completely disjoint as strings, three of five feeder slots contain domain-adjacent pairs (LOOM\,/\,FORGE from the technology domain; ONAGER\,/\,SIEGE from siege warfare; PATROL\,/\,MARCH from unit commands), suggesting that the content library exerts a partial anchoring force on word choice even when it does not constrain which word is chosen. Domain knowledge is neither necessary nor sufficient for stability; specification is both.

This paper makes five contributions. First, we introduce the same-input divergence paradigm as a general experimental method for characterizing the stability structure of any AI creative pipeline --- a complement to prompt sensitivity analysis (which varies the input to measure output sensitivity) that instead fixes the input and characterizes natural output variation across runs. Second, we provide concrete empirical data on stable and variable decisions at five levels of abstraction in a complete AI puzzle hunt pipeline, grounding the stability analysis in the specific numerical results of the AoE2 comparison: structural and assignment-level divergence is near-total on exact-string metrics; quality-level convergence is near-complete on pass/fail and normalized aggregate scores. Third, we establish the specification-drives-convergence finding as the primary explanatory principle: stability tracks what the brief said, not what the domain constrains, in this pipeline. Fourth, we characterize domain semantic gravity as a distinct phenomenon that operates below the level of specification but above the level of random variation, anchoring word-choice distributions without fixing individual words. Fifth, we derive actionable practical recommendations from the stability data: run once and trust for explicitly specified decisions; run multiple times and select for unspecified ones; use specification completeness as the primary lever for controlling pipeline reliability.
